\documentclass[11pt]{article}

% Shared notes settings for Elements of AIML lectures (UPES).
% Keep this file free of \documentclass to allow reuse via \input.

\usepackage[utf8]{inputenc}
\usepackage[T1]{fontenc}
\usepackage[a4paper,margin=1in]{geometry}
\usepackage{hyperref}
\usepackage{xurl}
\usepackage{booktabs}
\usepackage{amsmath}
\usepackage{amssymb}
\usepackage{listings}
\usepackage{xcolor}
\usepackage{enumitem}
\usepackage{parskip}

% Code listings (general-purpose).
\lstset{
  basicstyle=\ttfamily\small,
  keywordstyle=\color{blue},
  commentstyle=\color{gray},
  stringstyle=\color{teal},
  showstringspaces=false,
  columns=fullflexible,
  frame=single,
  framerule=0.2pt,
  breaklines=true
}

\setlist[itemize]{topsep=0.3em,itemsep=0.2em}
\setlist[enumerate]{topsep=0.3em,itemsep=0.2em}

% Simple per-section source line.
\newcommand{\notesources}[1]{\par\noindent{\small\color{gray}\textit{Sources:} \sloppy #1\par}}

% Unit-wise source strings for this subject (used by slides/notes).
% Path is relative to the lecture `latex/` folder (the compilation CWD).
\IfFileExists{../../../_shared/aiml-sources.tex}{%
  % Source strings for Elements of AIML (used by slides + notes).

\newcommand{\AIMLRepoURL}{https://github.com/tali7c/Elements-of-AIML}

\newcommand{\AIMLLabSources}{%
Provost and Fawcett, \textit{Data Science for Business}.;%
McKinney, \textit{Python for Data Analysis}.;%
WEKA Documentation (Waikato Environment for Knowledge Analysis), accessed 2026-02-28.;%
Matplotlib and Seaborn documentation, accessed 2026-02-28.%
}

\newcommand{\AIMLUnitISources}{%
Rich and Knight, \textit{Artificial Intelligence}, McGraw Hill, 2017.;%
Russell and Norvig, \textit{Artificial Intelligence: A Modern Approach} (4e), Ch. 1--2 (Introduction, Intelligent Agents).;%
Poole and Mackworth, \textit{Artificial Intelligence: Foundations of Computational Agents} (2e), selected sections.%
}

\newcommand{\AIMLUnitIISources}{%
Russell and Norvig, \textit{Artificial Intelligence: A Modern Approach} (4e), Ch. 7--9 (Logical Agents, First-Order Logic, Inference).;%
Huth and Ryan, \textit{Logic in Computer Science} (2e), selected sections on propositional and first-order logic.;%
Genesereth and Nilsson, \textit{Logical Foundations of Artificial Intelligence}, selected chapters.%
}

\newcommand{\AIMLUnitIIISources}{%
Mueller and Massaron, \textit{Machine Learning for Dummies}, For Dummies, 2016.;%
Mitchell, \textit{Machine Learning}, selected chapters on learning basics and evaluation.;%
Scikit-learn documentation, accessed 2026-02-28.%
}

\newcommand{\AIMLUnitIVSources}{%
Mueller and Massaron, \textit{Machine Learning for Dummies}, For Dummies, 2016.;%
James, Witten, Hastie, and Tibshirani, \textit{An Introduction to Statistical Learning} (2e), selected chapters.;%
Scikit-learn documentation, accessed 2026-02-28.%
}

\newcommand{\AIMLUnitVSources}{%
NIST, \textit{AI Risk Management Framework (AI RMF 1.0)}, 2023.;%
Barocas, Hardt, and Narayanan, \textit{Fairness and Machine Learning} (online book), selected sections.;%
Knaflic, \textit{Storytelling with Data}, selected chapters on communicating results.%
}
%
}{%
  % Faculty copies live under `private/` so their compile CWD is different.
  \IfFileExists{../../../../public/_shared/aiml-sources.tex}{%
    % Source strings for Elements of AIML (used by slides + notes).

\newcommand{\AIMLRepoURL}{https://github.com/tali7c/Elements-of-AIML}

\newcommand{\AIMLLabSources}{%
Provost and Fawcett, \textit{Data Science for Business}.;%
McKinney, \textit{Python for Data Analysis}.;%
WEKA Documentation (Waikato Environment for Knowledge Analysis), accessed 2026-02-28.;%
Matplotlib and Seaborn documentation, accessed 2026-02-28.%
}

\newcommand{\AIMLUnitISources}{%
Rich and Knight, \textit{Artificial Intelligence}, McGraw Hill, 2017.;%
Russell and Norvig, \textit{Artificial Intelligence: A Modern Approach} (4e), Ch. 1--2 (Introduction, Intelligent Agents).;%
Poole and Mackworth, \textit{Artificial Intelligence: Foundations of Computational Agents} (2e), selected sections.%
}

\newcommand{\AIMLUnitIISources}{%
Russell and Norvig, \textit{Artificial Intelligence: A Modern Approach} (4e), Ch. 7--9 (Logical Agents, First-Order Logic, Inference).;%
Huth and Ryan, \textit{Logic in Computer Science} (2e), selected sections on propositional and first-order logic.;%
Genesereth and Nilsson, \textit{Logical Foundations of Artificial Intelligence}, selected chapters.%
}

\newcommand{\AIMLUnitIIISources}{%
Mueller and Massaron, \textit{Machine Learning for Dummies}, For Dummies, 2016.;%
Mitchell, \textit{Machine Learning}, selected chapters on learning basics and evaluation.;%
Scikit-learn documentation, accessed 2026-02-28.%
}

\newcommand{\AIMLUnitIVSources}{%
Mueller and Massaron, \textit{Machine Learning for Dummies}, For Dummies, 2016.;%
James, Witten, Hastie, and Tibshirani, \textit{An Introduction to Statistical Learning} (2e), selected chapters.;%
Scikit-learn documentation, accessed 2026-02-28.%
}

\newcommand{\AIMLUnitVSources}{%
NIST, \textit{AI Risk Management Framework (AI RMF 1.0)}, 2023.;%
Barocas, Hardt, and Narayanan, \textit{Fairness and Machine Learning} (online book), selected sections.;%
Knaflic, \textit{Storytelling with Data}, selected chapters on communicating results.%
}
%
  }{%
    % Source strings for Elements of AIML (used by slides + notes).

\newcommand{\AIMLRepoURL}{https://github.com/tali7c/Elements-of-AIML}

\newcommand{\AIMLLabSources}{%
Provost and Fawcett, \textit{Data Science for Business}.;%
McKinney, \textit{Python for Data Analysis}.;%
WEKA Documentation (Waikato Environment for Knowledge Analysis), accessed 2026-02-28.;%
Matplotlib and Seaborn documentation, accessed 2026-02-28.%
}

\newcommand{\AIMLUnitISources}{%
Rich and Knight, \textit{Artificial Intelligence}, McGraw Hill, 2017.;%
Russell and Norvig, \textit{Artificial Intelligence: A Modern Approach} (4e), Ch. 1--2 (Introduction, Intelligent Agents).;%
Poole and Mackworth, \textit{Artificial Intelligence: Foundations of Computational Agents} (2e), selected sections.%
}

\newcommand{\AIMLUnitIISources}{%
Russell and Norvig, \textit{Artificial Intelligence: A Modern Approach} (4e), Ch. 7--9 (Logical Agents, First-Order Logic, Inference).;%
Huth and Ryan, \textit{Logic in Computer Science} (2e), selected sections on propositional and first-order logic.;%
Genesereth and Nilsson, \textit{Logical Foundations of Artificial Intelligence}, selected chapters.%
}

\newcommand{\AIMLUnitIIISources}{%
Mueller and Massaron, \textit{Machine Learning for Dummies}, For Dummies, 2016.;%
Mitchell, \textit{Machine Learning}, selected chapters on learning basics and evaluation.;%
Scikit-learn documentation, accessed 2026-02-28.%
}

\newcommand{\AIMLUnitIVSources}{%
Mueller and Massaron, \textit{Machine Learning for Dummies}, For Dummies, 2016.;%
James, Witten, Hastie, and Tibshirani, \textit{An Introduction to Statistical Learning} (2e), selected chapters.;%
Scikit-learn documentation, accessed 2026-02-28.%
}

\newcommand{\AIMLUnitVSources}{%
NIST, \textit{AI Risk Management Framework (AI RMF 1.0)}, 2023.;%
Barocas, Hardt, and Narayanan, \textit{Fairness and Machine Learning} (online book), selected sections.;%
Knaflic, \textit{Storytelling with Data}, selected chapters on communicating results.%
}
%
  }%
}


\title{Elements of AIML\\Unit 05 -- Lecture 03 Notes\\Applications in Education, Insurance, Retail, Transport \& Energy (Wrap-up)}
\author{Tofik Ali}
\date{\today}

\begin{document}
\maketitle
\tableofcontents

\section{Lecture Overview}
This final lecture completes Unit V by surveying additional application domains and providing:
\begin{itemize}[leftmargin=*]
  \item a practical deployment checklist (what can go wrong),
  \item a revision set that connects Units I--V.
\end{itemize}
\notesources{\AIMLUnitVSources}

\section{How to use these notes (self-study)}
Use this lecture as an integration lesson.
Pick one domain (education/insurance/retail/transport/energy) and apply the checklist: define the task, pick metrics, identify risks, and propose a monitoring plan.
Then use the revision questions to connect logic (Unit II) with evaluation (Unit III) and modeling (Unit IV).

\section{Learning Outcomes}
\begin{itemize}[leftmargin=*]
  \item List use-cases in education, insurance, retail/supply chain, transport, and energy.
  \item Explain domain constraints (metrics, drift, safety, privacy).
  \item Apply a checklist to evaluate deployment risks.
  \item Answer revision questions across the course.
\end{itemize}

\section{Before you start}
Treat this lecture as a course-integration pack:
\begin{itemize}[leftmargin=*]
  \item from \textbf{Unit 05 Lecture 01}, reuse the fairness/privacy/accountability lens,
  \item from \textbf{Unit 05 Lecture 02}, reuse the domain-case template of task, metric, risk, and monitoring,
  \item from Units II--IV, keep the underlying reasoning, evaluation, and modeling tools available for the revision set.
\end{itemize}

\section{Keywords}
\begin{itemize}[leftmargin=*]
  \item \textbf{Drift:} changing patterns over time.
  \item \textbf{Feedback loop:} model outputs influence future data (recommendations).
  \item \textbf{Checklist:} structured risk review for deployment.
\end{itemize}

\section{Abbreviations and key terms}
\begin{itemize}[leftmargin=*]
  \item \textbf{KPI}: key performance indicator (business success measure).
  \item \textbf{ETA}: estimated time of arrival (transport).
  \item \textbf{SLA}: service-level agreement (latency/uptime requirements).
  \item \textbf{Objective mismatch}: optimizing the wrong metric for the real-world goal.
\end{itemize}

\section{Cross-domain case map}
Use the same lens across every domain:
\begin{itemize}[leftmargin=*]
  \item \textbf{Education:} task = recommendation or dropout prediction; metric = completion/engagement plus fairness; risk = privacy and unfair intervention; monitoring = subgroup performance and engagement trends.
  \item \textbf{Insurance:} task = pricing support, fraud detection, or claim triage; metric = cost-aware classification/regression plus explainability; risk = proxy bias and regulatory failure; monitoring = appeal rates and subgroup gaps.
  \item \textbf{Retail/supply chain:} task = forecasting, recommendation, or inventory optimization; metric = MAE/RMSE or ranking/business KPIs; risk = seasonality, feedback loops, and data quality; monitoring = forecast error and product-exposure drift.
  \item \textbf{Transportation/logistics:} task = ETA prediction, routing, or predictive maintenance; metric = ETA MAE/RMSE and SLA compliance; risk = sensor failure and safety issues; monitoring = latency, route-error trends, and sensor health.
  \item \textbf{Energy/utilities:} task = load forecasting or fault detection; metric = forecasting error or recall on rare events; risk = reliability failure under rare conditions; monitoring = demand-pattern drift and alarm quality.
\end{itemize}

\section{Applications across domains}
\subsection{Education}
\begin{itemize}[leftmargin=*]
  \item personalization and recommendation,
  \item early warning for dropout risk,
  \item analysis of learning outcomes.
\end{itemize}
Important constraints: privacy, fairness, transparency.

\paragraph{Typical data/labels (examples).}
Learning activity logs, assessment scores, attendance, and engagement signals.
Labels could be ``dropout'' or ``course completion''.
Metrics must consider fairness (avoid disadvantaging groups) and privacy of student data.

\subsection{Insurance}
\begin{itemize}[leftmargin=*]
  \item risk modeling and pricing support,
  \item claim triage and fraud detection,
  \item customer support automation.
\end{itemize}
Important constraints: explainability, fairness, regulation.

\paragraph{Typical tasks and risks.}
Fraud detection is imbalanced; pricing models require transparency and governance.
Proxy features can create unfair outcomes, and models must support audit/appeals.

\subsection{Retail and supply chain}
\begin{itemize}[leftmargin=*]
  \item demand forecasting and inventory management,
  \item recommendation systems,
  \item logistics optimization.
\end{itemize}
Important constraints: seasonality/drift, feedback loops, data quality.

\paragraph{Feedback loop example.}
If recommendations increase exposure of some products, the data collected reflects the system's own choices, which can amplify popularity bias unless corrected.

\subsection{Transportation and logistics}
\begin{itemize}[leftmargin=*]
  \item ETA prediction, routing,
  \item predictive maintenance,
  \item safety monitoring.
\end{itemize}
Important constraints: real-time requirements, safety, sensor errors.

\paragraph{Metrics.}
ETA uses MAE/RMSE (minutes). Routing/dispatch may optimize cost/time subject to safety and SLAs.
Sensor failures and drift require monitoring and fallback strategies.

\subsection{Energy and utilities}
\begin{itemize}[leftmargin=*]
  \item load forecasting,
  \item fault detection,
  \item demand-response optimization.
\end{itemize}
Important constraints: reliability, rare events, monitoring.

\paragraph{Rare-event issue.}
Fault detection resembles imbalanced classification/anomaly detection: most time is normal operation, so accuracy can be misleading.

\section{What could go wrong? (deployment checklist)}
Use this checklist before and after deployment:
\begin{enumerate}[leftmargin=*]
  \item \textbf{Objective mismatch:} are we optimizing the right metric/cost?
  \item \textbf{Data problems:} leakage, bias, missing values, non-representative data.
  \item \textbf{Fairness/privacy:} are there subgroup gaps? are we compliant?
  \item \textbf{Drift:} what changes over time? what is the monitoring plan?
  \item \textbf{Security/misuse:} adversarial inputs, abuse, unauthorized access.
\end{enumerate}

\section{Common pitfalls (course wrap-up)}
\begin{itemize}[leftmargin=*]
  \item \textbf{Thinking ``accuracy is everything'':} real systems require the right metric, fairness checks, and cost-aware decisions.
  \item \textbf{Ignoring feedback loops:} recommendations and policies change user behavior, which changes the future data.
  \item \textbf{No monitoring/retraining plan:} drift is normal; lack of monitoring turns small issues into failures.
  \item \textbf{Skipping documentation:} without documentation, auditing and debugging deployed models becomes very difficult.
\end{itemize}

\section{Practice (with answers)}
\subsection*{P1. Give one KPI for an education recommender system.}
\textbf{Answer:} improvement in course completion rate, or increased engagement time, while ensuring fairness across groups.

\subsection*{P2. Why are time-based splits important for forecasting tasks?}
\textbf{Answer:} the model will be used on future data; random shuffling can leak future patterns into training and produce overly optimistic estimates.

\section{Revision questions (with answers)}

\subsection*{R1. What is entailment and how can we test it using refutation?}
\textbf{Answer:} entailment $KB \models \alpha$ means $\alpha$ is true in every model of $KB$. Refutation test: check if $KB \land \lnot \alpha$ is unsatisfiable (derive contradiction via resolution).
\textbf{Revisit:} Unit 02 Lecture 03.

\subsection*{R2. Why do we need train/validation/test splits?}
\textbf{Answer:} to measure generalization honestly. Train fits parameters, validation tunes choices, test gives an unbiased final evaluation. Reusing test for tuning causes optimistic bias.
\textbf{Revisit:} Unit 03 Lecture 02.

\subsection*{R3. Regression vs classification (and metrics)?}
\textbf{Answer:} regression predicts continuous values (MAE/RMSE/$R^2$). Classification predicts labels (precision/recall/F1/ROC-AUC). Metric choice depends on error costs and imbalance.
\textbf{Revisit:} Unit 04 Lectures 01--02.

\subsection*{R4. What does PCA optimize and why standardize first?}
\textbf{Answer:} PCA finds orthogonal directions of maximum variance. Standardization prevents large-scale features from dominating the covariance and components.
\textbf{Revisit:} Unit 03 Lecture 03.

\subsection*{R5. Give two examples of drift and one monitoring method.}
\textbf{Answer:} fraud patterns change; customer behavior changes seasonally; spam tactics evolve. Monitoring: track feature distributions and performance metrics over time with alerts.
\textbf{Revisit:} Unit 03 Lecture 02 and Unit 05 Lecture 02.

\subsection*{R6. What is class imbalance and why is accuracy misleading?}
\textbf{Answer:} imbalance means one class is rare. Accuracy can be high by predicting the majority class; use precision/recall/F1 and cost-aware evaluation.
\textbf{Revisit:} Unit 04 Lecture 02.

\subsection*{R7. What is unification used for in logic inference?}
\textbf{Answer:} unification finds substitutions to match FOL literals so resolution can be applied with variables (e.g., $Human(x)$ with $Human(Socrates)$).
\textbf{Revisit:} Unit 02 Lecture 03.

\subsection*{R8. Give one example of a feedback loop in recommendations.}
\textbf{Answer:} if a recommender shows certain products more, users click them more, which produces data that further reinforces recommending the same products, potentially reducing diversity and increasing bias.
\textbf{Revisit:} Unit 05 Lecture 03, Retail and Supply Chain.

\subsection*{R9. Why is explainability important in insurance?}
\textbf{Answer:} decisions affect people financially; regulations and trust require understandable reasons; explainability supports auditing and contesting errors.
\textbf{Revisit:} Unit 05 Lecture 03, Insurance.

\subsection*{R10. What are two responsible AI checks before deployment?}
\textbf{Answer:} fairness evaluation across subgroups and privacy/security review (data governance, access control), plus monitoring plan for drift.
\textbf{Revisit:} Unit 05 Lecture 01.

\section{Exit questions (with answers)}

\subsection*{Q1. Pick one domain and propose an ML task + suitable metric.}
\textbf{Answer:} retail demand forecasting (regression) with RMSE/MAE; fraud detection (classification) with precision/recall; healthcare screening with high recall.

\subsection*{Q2. List three checklist items.}
\textbf{Answer:} objective mismatch, data leakage/bias, drift and monitoring, fairness/privacy, security/misuse (any three).

\subsection*{Q3. Why explainability in insurance?}
\textbf{Answer:} accountability and regulation; decisions must be justified and audited.

\subsection*{Q4. One example of feedback loop in retail recommendations.}
\textbf{Answer:} popular items get more exposure, then more clicks, then even more exposure, causing self-reinforcing popularity bias.

\section{Summary}
This course combined logic-based reasoning and data-driven learning, with emphasis on evaluation and responsible application in real domains.
\notesources{\AIMLUnitVSources}

\end{document}
